% ---
% id: meat-core/frank-wolfe/rem:fo-iterate-checks
% kind: rem
% title: "What the iterate checks establish"
% summary: >-
%   The three iterate checks decide every condition of the relaxed feasible
%   set, all of them on the J factor; the acceptance factor holds by
%   construction. They serve the membership hypothesis and the estimate's
%   premise; no check of a solver's return would serve the bound.
% prominence: standard
% status: accepted
% uses: [thm:firstorder, def:first-order-setting, def:relaxed-feasible-set, meat-core/thm:choi, prop:relaxed-set-properties, rem:acc-side]
% mentions: [req:iterate-checks, req:iterate-repair, req:acceptance-iterate-exact, req:verified-feasible-side, foundations/req:tolerance-scope, rem:fo-returned-point]
% assumes: []
% refs: []
% provenance: newly-derived
% lean:
% ---

\begin{remark}\label{rem:fo-iterate-checks}
Requirement~\ref{req:iterate-checks} decides, on the point consumed at each
iterate, every condition defining $\Feas^{J}$ in \eqref{eq:choi-feas}:
per-block positive semidefiniteness, the coupling condition
$\sum_s\tr_{\Qsys_s}J_s=\idon{A_0}$, \emph{and the declared constraint rows}
$\bowtie_k$. Checking two of the three is not checking the condition.
\emph{No such check is made on a point a solver returned, and none would
serve the bound}: verifying $S_t\in\Feas^{J}$ gives
$\langle\lmoJ_t,S_t\rangle\ge\min_{J'\in\Feas^{J}}\langle\lmoJ_t,J'\rangle$,
the \emph{opposite} inequality to \eqref{eq:lmo}. Feasibility bounds the
linear form from below, which is exactly the side the first-order claim
cannot use; one can verify \emph{which set} a solver minimised over, never
that it minimised. The feasible side of a bracket is therefore a lower bound
formed without any point of the solver
(Requirement~\ref{req:verified-feasible-side}). What the checks do buy is
the membership hypothesis of Theorem~\ref{thm:firstorder}\,(iiib) on its
$J$ factor, and Theorem~\ref{thm:firstorder}\,(i),
respectively~(iiia): the point the run reports its estimate at is a convex
combination of oracle returns; \eqref{eq:fo-estimate} holds only for points
of $\acSet\times\Feas^{J}$, which is what a run does \emph{not} establish,
and (iiia) is the part it is in, giving
$O^{J}_{\alpha}(x)\ge\hlow[\Feas^{J}(\varepsilon)]$ at the relaxed
membership the checks do establish --- on the point the run returned, the
membership the end-of-run measurement of
Remark~\ref{rem:fo-returned-point} establishes, which is of the $J$ factor
only. So the checks serve the bound's membership hypothesis and the premise
of the estimate tier, which is the tier reported on every run. Their windows are numbers of their own and not
the declared tolerance $\tnum$, the object checked being one the optimiser
computed (\xref{foundations/req:tolerance-scope}\,(ii)).

Between them the three checks are what membership in
$\Feas^{J}(\varepsilon)$ at the larger of the two declared windows asks
for, and they cover all of it: per-block positive semidefiniteness, which
the family does not relax and which is decided on the point consumed at
exactly $0$ and at no window (Requirement~\ref{req:iterate-repair} says
what may be done to a block that fails it, and the record it must carry);
the coupling residual and each $\bowtie_k$ residual, each inside its own
window, hence inside the larger of the two by monotonicity
(Proposition~\ref{prop:relaxed-set-properties}\,(a)); and the family's
normalisation clause with them, that clause being implied by the coupling
clause and adding no check
(Proposition~\ref{prop:relaxed-set-properties}\,(c$'$)). No fourth check is
needed.

\emph{All three are conditions on the $J$ factor.} The hypothesis they
serve is on a \emph{product}, $x_t\in\acSet\times\Feas^{J}(\varepsilon)$,
and nothing in these checks touches the other factor: \emph{the three
checks establish the $J$ half and no more, and the $\acSet$ half is not
checked at all}. It is instead true \emph{by construction}, the
acceptance-side iterate being formed in exact rational arithmetic on the
\tref{exact island} under ($\alpha$) and ($\beta$) of
Remark~\ref{rem:acc-side}: a convex combination, with exact rational
weights, of exact points of a convex set, so $\mathbf q_t\in\acSet$ by
induction. Remark~\ref{rem:acc-side} names the three consumers of
Theorem~\ref{thm:firstorder} that break if it does not hold. A run that
forms $\mathbf q_t$ in binary64 instead has not met the membership
hypothesis (Requirement~\ref{req:acceptance-iterate-exact}); there is no
check for it to pass on that factor, and none is offered. The coverage
claim of the previous paragraph is therefore a claim about one factor of a
product.

\emph{The window on the positivity check decides whether the coverage claim
is true.} At a positive window $w$ the positivity check would accept a
point with $\lmin\in(-w,0)$, which lies in \emph{no}
$\Feas^{J}(\varepsilon)$ for any $\varepsilon$ --- the family relaxes the
coupling and the $\bowtie_k$ rows and relaxes positivity not at all
(Proposition~\ref{prop:relaxed-set-properties}\,(c), \emph{Fourth}) --- so
at such a point $O^{J}_{\alpha}$ is not even defined and neither (i) nor
(ii) nor (iii) of Theorem~\ref{thm:firstorder} applies. That is why
positivity is decided at $0$, and the drift a declared window would let
pass silently is repaired and recorded under
Requirement~\ref{req:iterate-repair}. By
Theorem~\ref{thm:firstorder}\,(iiia) that reading is the conservative one:
membership in the wider $\Feas^{J}(\varepsilon)$ gives
$O^{J}_{\alpha}(x)\ge\hlow[\Feas^{J}(\varepsilon)]$, the \emph{Choi-form}
relaxed infimum \eqref{eq:hlow-feasJ}, and
$\hlow[\Feas^{J}(\varepsilon)]\le\hlow$ by \eqref{eq:feas-tol-mono-J}, so
the estimate over-estimates a quantity that is itself no larger than
$\hlow$. (The state-form $\hlow[\Feas(\varepsilon)]$ is a different number
and is not what the checks buy;
Proposition~\ref{prop:relaxed-set-properties}\,(d$'$) says why the two are
not interchangeable at $\varepsilon>0$.)
\end{remark}
